\chapter{Введение}

Современное компьютерное зрение уже давно не сводится только к подбору архитектуры с максимальной точностью на тестовой выборке. Да, итоговое качество модели остаётся главным практическим показателем. Но оно почти ничего не говорит о том, как именно модель преобразует изображение внутри себя. Две архитектуры могут давать близкий ответ на выходе, но строить совершенно разные промежуточные представления. Собственно, это и делает анализ внутренних признаков отдельной исследовательской задачей.

За последние годы в задачах зрения сложилось несколько крупных архитектурных семейств. Сначала основным инструментом были свёрточные нейронные сети. Затем широкое распространение получили остаточные сети ResNet, которые позволили устойчиво обучать более глубокие CNN за счёт остаточных связей \cite{he2016deep}. Позже в компьютерное зрение пришли визуальные трансформеры. Vision Transformer разбивает изображение на патчи и обрабатывает их как последовательность токенов \cite{dosovitskiy2021image}; эта идея опирается на архитектуру Transformer, предложенную для последовательностей \cite{vaswani2017attention}. Ещё один важный вариант — Swin Transformer, где трансформерная обработка сочетается с иерархией, оконным вниманием, сдвигом окон и объединением патчей \cite{liu2021swin}. Получается три разных способа работы с изображением: свёрточная пирамида, токенная модель с фиксированной решёткой патчей и иерархический трансформер.

Внутри нейронной сети изображение постепенно превращается в набор признаков. В ранних слоях эти признаки часто ближе к локальным изменениям: границам, текстурам, цветовым переходам. На более глубоких уровнях представления обычно становятся более агрегированными. Но здесь нельзя делать слишком прямой вывод. Глубокий слой — это не обязательно «только семантика», а ранний — не обязательно «только текстура». Это скорее удобная интуиция, которую нужно проверять количественно.

Один из способов такой проверки — перейти к частотной области. Изображение можно рассматривать как сигнал, состоящий из компонент разного масштаба. Низкие пространственные частоты связаны с плавными изменениями и крупными структурами, более высокие — с резкими переходами, границами, текстурами и шумовыми компонентами. Если рассматривать не только исходное изображение, но и внутренние представления модели, появляется естественный вопрос: как по глубине меняется распределение энергии по частотам?

Для CNN такой вопрос возникает почти сразу. Карты признаков имеют двумерную пространственную структуру, а сама операция свёртки тесно связана с фильтрацией сигнала. При этом в ResNet между стадиями уменьшается пространственное разрешение. Значит, если в глубоких слоях виден более узкий спектр, нельзя автоматически считать это только свойством обученных признаков. Часть эффекта может быть обычным следствием уменьшения размера карты признаков. Поэтому в данной работе отдельно проверяется, насколько спектральное сжатие ResNet сохраняется после контроля этого фактора.

С ViT ситуация другая. Модель работает с последовательностью патчевых токенов и в стандартной конфигурации сохраняет одну и ту же патчевую решётку по глубине. Это делает ViT удобным контрастом к ResNet: если спектральный профиль меняется между блоками, то это изменение уже нельзя объяснить межстадийным уменьшением пространственного разрешения. При этом не стоит сразу приписывать наблюдаемую динамику одному только самовниманию. На результат влияют остаточные связи, нормализации, MLP-блоки, позиционные эмбеддинги и способ кодирования патчей. Работы по визуальным трансформерам уже показывали, что их внутренние представления отличаются от CNN \cite{raghu2021do}, а механизм самовнимания может проявлять низкочастотные свойства \cite{park2022how}. Но это не означает, что все блоки ViT обязаны монотонно смещаться к низким частотам.

Swin Transformer занимает промежуточное положение. Он использует токены и механизм внимания, но при этом строит иерархию и уменьшает размер токенной решётки через объединение патчей \cite{liu2021swin}. Поэтому для Swin особенно важно не смешивать два эффекта: изменение самих признаков и уменьшение пространственного разрешения. В этой работе Swin рассматривается именно как архитектура, где эти два фактора нужно разделять аккуратно.

Теоретическая мотивация исследования связана с работами о частотном смещении нейронных сетей. В \cite{rahaman2019spectral} показано, что нейронные сети обычно быстрее осваивают низкочастотные компоненты функций. Позже эта линия развивалась в исследованиях принципа частотного обучения и способов измерения частотного смещения \cite{xu2025overview,kiessling2022computable}. В задачах компьютерного зрения частотные компоненты также связывают с устойчивостью и обобщающей способностью моделей \cite{yin2019fourier,wang2020highfrequency}. В данной работе эти результаты используются как основание для постановки вопроса, но не как готовое объяснение. Спектр показывает, как распределена энергия. Он не доказывает сам по себе, какие частоты полезны для классификации.

Отсюда и появляется основная идея работы. Нужно не просто построить красивые спектральные кривые, а сделать сравнение доказательным: использовать единый протокол для разных архитектур, корректно считать радиальные метрики, учитывать число точек в радиальных кольцах, задавать единый частотный масштаб и проверять влияние уменьшения разрешения там, где оно есть. Без этого легко перепутать реальное изменение представлений с техническим эффектом пространственной дискретизации.

Цель данной выпускной квалификационной работы — исследовать спектральные характеристики внутренних представлений ResNet, ViT и Swin и определить, как тип архитектуры связан с перераспределением спектральной энергии по глубине модели.

Для достижения этой цели были поставлены следующие задачи:

\begin{enumerate}
\item изучить теоретические основы спектрального анализа нейросетевых представлений и связанные исследования по частотному смещению, CNN и визуальным трансформерам;
\item описать архитектурные особенности ResNet, ViT и Swin, влияющие на пространственную структуру представлений;
\item реализовать единый протокол извлечения промежуточных представлений из моделей разных архитектур;
\item адаптировать спектральный анализ к токенным представлениям ViT и Swin через восстановление двумерной решётки патчей;
\item построить радиальные спектральные кривые для выбранных уровней моделей;
\item рассчитать спектральный центроид, долю низкочастотной энергии и коэффициент экспоненциального спектрального сжатия;
\item добавить стандартные отклонения по изображениям выборки и использовать их при интерпретации таблиц и графиков;
\item провести контрольные эксперименты для ResNet и Swin, чтобы отделить изменение спектрального профиля от простого уменьшения пространственного разрешения;
\item сравнить спектральную динамику внутри каждого архитектурного семейства и между семействами ResNet, ViT и Swin.
\end{enumerate}

Объектом исследования являются внутренние представления моделей компьютерного зрения. Предметом исследования являются спектральные характеристики этих представлений: распределение энергии по пространственным частотам, изменение радиального спектрального профиля по глубине модели и различия между архитектурными семействами.

В качестве экспериментальных объектов используются три семейства моделей. ResNet рассматривается как представитель иерархических свёрточных архитектур. Vision Transformer используется как пример модели, работающей с последовательностью патчевых токенов. Swin Transformer выбран как иерархический трансформер с локальной оконной обработкой и уменьшением токенной решётки. Эксперименты проводятся на подвыборке ImageNet 10K \cite{priyerana_imagenet10k}, основанной на классическом наборе ImageNet \cite{deng2009imagenet}.

В работе проверяются три исследовательских предположения.

Первое предположение состоит в том, что ResNet по мере перехода к глубоким стадиям демонстрирует снижение спектрального центроида и рост доли низкочастотной энергии. При этом отдельно проверяется, какая часть этой динамики связана с уменьшением пространственного разрешения между стадиями.

Второе предположение связано с ViT. Поскольку стандартный ViT сохраняет фиксированную решётку патчей по глубине, его спектральная динамика не обязана повторять монотонное низкочастотное сжатие CNN. Ожидается, что изменение спектра будет более немонотонным.

Третье предположение относится к Swin. Из-за иерархической структуры и объединения патчей Swin может приходить к более низкочастотному финальному профилю, но часть этого эффекта, возможно, будет объясняться уменьшением токенной решётки. Поэтому для Swin особенно нужен контроль влияния разрешения.

Научная новизна работы состоит в том, что спектральная динамика ResNet, ViT и Swin сравнивается в едином экспериментальном протоколе, но с учётом архитектурных различий. Для графиков используется радиальная средняя мощность, а для численных метрик — суммарная энергия по радиальным кольцам. Это снимает методическую проблему, при которой радиальное среднее ошибочно трактуется как доля энергии. Также в работе добавлены контрольные эксперименты для ResNet и Swin, позволяющие отделить эффект уменьшения разрешения от изменения самих представлений.

Практическая значимость работы состоит в том, что предложенный протокол можно использовать как дополнительный инструмент диагностики моделей компьютерного зрения. Он показывает, где по глубине модели происходит сдвиг энергии к низким частотам, где сохраняется среднечастотный хвост и насколько этот эффект связан с архитектурной иерархией. Такой анализ может быть полезен при сравнении архитектур для задач, где пространственная структура признаков имеет значение: медицинские изображения, спутниковые снимки, дефектоскопия, текстурный анализ и задачи с мелкими объектами. При этом в данной работе не делается прямой вывод о полезности конкретных частот для качества классификации. Для этого нужен отдельный эксперимент.

Структура работы следующая. Во второй главе вводятся теоретические сведения: внутренние представления, частотное описание изображений, радиальные спектры, круг Найквиста, ResNet, ViT, Swin и связанные исследования. В третьей главе описывается методика: извлечение представлений, построение спектра мощности, радиальное представление, расчёт метрик и контрольные эксперименты. В четвёртой главе приведены результаты: спектральные кривые, таблицы метрик, сравнение архитектур и проверка влияния уменьшения разрешения. В пятой главе формулируются основные выводы, ограничения и направления дальнейшей работы.
